% GENERATED FILE. Edit canonical lessons, then run npm run generate:books.
% canonical-source-hash: fnv1a64:54a618353b064f48
% canonical-lessons: PA-C81-saukha, PA-C81-mahinga, PA-C81-sasta, PA-C81-uchcha, PA-C81-nivam

\chapter{Easy, Expensive, Cheap, High, Low}
\label{ch:pa-easy-expensive-cheap-high-low}
\hlchaptermodality{81}

\begin{chapteropening}
\textbf{By the end of this chapter:} \emph{I can say easy, expensive, cheap, high and low.}

Complete the last lesson of chapter 81: I can say easy, expensive, cheap, high and low.
\end{chapteropening}

\section[\texorpdfstring{\pa{ਸੌਖਾ} (saukhā)}{ਸੌਖਾ (saukhā)}]{\pa{ਸੌਖਾ} (saukhā) — easy}

\label{lesson:PA-C81-saukha}



\begin{quote}
Before the new one: say the Punjabi for slowly, then the Punjabi for fast.
\end{quote}

\subsection*{You'll want to know: \pa{ਸੌਖਾ}}
\textbf{\pa{ਸੌਖਾ}} — \emph{saukhā} — \textquotedblleft{}easy\textquotedblright{}.

\begin{cousinweb}
Where \textbf{\pa{ਸੌਖਾ}} (\emph{saukhā}) comes from is not settled. CDIAL has no entry for it, though it lists Punjabi \emph{suhelā}, \textquotedblleft{}easy\textquotedblright{}, under Sanskrit \emph{sukha}, \textquotedblleft{}easy, pleasant\textquotedblright{} (CDIAL 13451).
\end{cousinweb}

\begin{grammarlens}[title={What you've built}]
One more word for this chapter.
\end{grammarlens}

\subsection*{Your turn}
\begin{itemize}
\raggedright
  \item \emph{Say it:} \emph{saukhā}
  \item \emph{Say it:} \emph{saukhā}, once more
  \item \emph{Say it:} say \emph{saukhā}
  \item \emph{Recall:} say the Punjabi for slowly, then the Punjabi for fast, then say \emph{saukhā} again
\end{itemize}

\begin{tcolorbox}[breakable,colback=teal!4,colframe=teal!35!black,title={Before you move on}]
What does \pa{ਸੌਖਾ} mean? (\textquotedblleft{}Easy\textquotedblright{}.) Say it once more.
\end{tcolorbox}
\section[\texorpdfstring{\pa{ਮਹਿੰਗਾ} (mahiṅgā)}{ਮਹਿੰਗਾ (mahiṅgā)}]{\pa{ਮਹਿੰਗਾ} (mahiṅgā) — expensive}

\label{lesson:PA-C81-mahinga}



\begin{quote}
Before the new one: say the Punjabi for fast, then the Punjabi for easy.
\end{quote}

\subsection*{You'll want to know: \pa{ਮਹਿੰਗਾ}}
\textbf{\pa{ਮਹਿੰਗਾ}} — \emph{mahiṅgā} — \textquotedblleft{}expensive\textquotedblright{}.

\begin{cousinweb}
\textbf{\pa{ਮਹਿੰਗਾ}} (\emph{mahiṅgā}) is a tadbhava, a word inherited from Sanskrit and reshaped by sound change. It comes from \emph{mahārgha}, \textquotedblleft{}high-priced\textquotedblright{} (CDIAL 9954). CDIAL lists Punjabi \emph{maĩhgā}.
\end{cousinweb}

\begin{grammarlens}[title={What you've built}]
One more word for this chapter.
\end{grammarlens}

\subsection*{Your turn}
\begin{itemize}
\raggedright
  \item \emph{Say it:} \emph{mahiṅgā}
  \item \emph{Say it:} \emph{mahiṅgā}, once more
  \item \emph{Say it:} say \emph{mahiṅgā}
  \item \emph{Recall:} say the Punjabi for fast, then the Punjabi for easy, then say \emph{mahiṅgā} again
\end{itemize}

\begin{tcolorbox}[breakable,colback=teal!4,colframe=teal!35!black,title={Before you move on}]
What does \pa{ਮਹਿੰਗਾ} mean? (\textquotedblleft{}Expensive\textquotedblright{}.) Say it once more.
\end{tcolorbox}
\section[\texorpdfstring{\pa{ਸਸਤਾ} (sastā)}{ਸਸਤਾ (sastā)}]{\pa{ਸਸਤਾ} (sastā) — cheap}

\label{lesson:PA-C81-sasta}



\begin{quote}
Before the new one: say the Punjabi for easy, then the Punjabi for expensive.
\end{quote}

\subsection*{You'll want to know: \pa{ਸਸਤਾ}}
\textbf{\pa{ਸਸਤਾ}} — \emph{sastā} — \textquotedblleft{}cheap\textquotedblright{}.

\begin{cousinweb}
Where \textbf{\pa{ਸਸਤਾ}} (\emph{sastā}) comes from is not settled: CDIAL has no entry for it.
\end{cousinweb}

\begin{grammarlens}[title={What you've built}]
One more word for this chapter.
\end{grammarlens}

\subsection*{Your turn}
\begin{itemize}
\raggedright
  \item \emph{Say it:} \emph{sastā}
  \item \emph{Say it:} \emph{sastā}, once more
  \item \emph{Say it:} say \emph{sastā}
  \item \emph{Recall:} say the Punjabi for easy, then the Punjabi for expensive, then say \emph{sastā} again
\end{itemize}

\begin{tcolorbox}[breakable,colback=teal!4,colframe=teal!35!black,title={Before you move on}]
What does \pa{ਸਸਤਾ} mean? (\textquotedblleft{}Cheap\textquotedblright{}.) Say it once more.
\end{tcolorbox}
\section[\texorpdfstring{\pa{ਉੱਚਾ} (uchchā)}{ਉੱਚਾ (uchchā)}]{\pa{ਉੱਚਾ} (uchchā) — high}

\label{lesson:PA-C81-uchcha}



\begin{quote}
Before the new one: say the Punjabi for expensive, then the Punjabi for cheap.
\end{quote}

\subsection*{You'll want to know: \pa{ਉੱਚਾ}}
\textbf{\pa{ਉੱਚਾ}} — \emph{uchchā} — \textquotedblleft{}high\textquotedblright{}.

\begin{cousinweb}
\textbf{\pa{ਉੱਚਾ}} (\emph{uchchā}) is a tadbhava, a word inherited from Sanskrit and reshaped by sound change. It comes from \emph{ucca}, \textquotedblleft{}tall, lofty\textquotedblright{} (CDIAL 1634).
\end{cousinweb}

\begin{grammarlens}[title={What you've built}]
One more word for this chapter.
\end{grammarlens}

\subsection*{Your turn}
\begin{itemize}
\raggedright
  \item \emph{Say it:} \emph{uchchā}
  \item \emph{Say it:} \emph{uchchā}, once more
  \item \emph{Say it:} say \emph{uchchā}
  \item \emph{Recall:} say the Punjabi for expensive, then the Punjabi for cheap, then say \emph{uchchā} again
\end{itemize}

\begin{tcolorbox}[breakable,colback=teal!4,colframe=teal!35!black,title={Before you move on}]
What does \pa{ਉੱਚਾ} mean? (\textquotedblleft{}High\textquotedblright{}.) Say it once more.
\end{tcolorbox}
\section[\texorpdfstring{\pa{ਨੀਵਾਂ} (nīvāṁ)}{ਨੀਵਾਂ (nīvāṁ)}]{\pa{ਨੀਵਾਂ} (nīvāṁ) — low}

\label{lesson:PA-C81-nivam}



\begin{quote}
Before the new one: say the Punjabi for cheap, then the Punjabi for high.
\end{quote}

\subsection*{You'll want to know: \pa{ਨੀਵਾਂ}}
\textbf{\pa{ਨੀਵਾਂ}} — \emph{nīvāṁ} — \textquotedblleft{}low\textquotedblright{}.

\begin{cousinweb}
\textbf{\pa{ਨੀਵਾਂ}} (\emph{nīvāṁ}) is a tadbhava, a word inherited from Sanskrit and reshaped by sound change. It comes from \emph{nīpa}, \textquotedblleft{}situated low, deep\textquotedblright{} (CDIAL 7548).
\end{cousinweb}

\begin{grammarlens}[title={What you've built}]
One more word for this chapter.
\end{grammarlens}

\subsection*{Your turn}
\begin{itemize}
\raggedright
  \item \emph{Say it:} \emph{nīvāṁ}
  \item \emph{Say it:} \emph{nīvāṁ}, once more
  \item \emph{Say it:} say \emph{nīvāṁ}
  \item \emph{Recall:} say the Punjabi for cheap, then the Punjabi for high, then say \emph{nīvāṁ} again
\end{itemize}

\begin{tcolorbox}[breakable,colback=teal!4,colframe=teal!35!black,title={Before you move on}]
What does \pa{ਨੀਵਾਂ} mean? (\textquotedblleft{}Low\textquotedblright{}.) Say it once more.
\end{tcolorbox}
